% !TeX program = xelatex
%==============================================================
%  Python金融数据分析 · 第4章 pandas数据分析
%  基于《Python金融大数据分析（第2版）》第5章内容改编
%  授课对象：本科金融系学生（已学完第2、3章）
%  编译：xelatex chapter4.tex（编译两遍）
%==============================================================
\documentclass[9pt]{ctexbeamer}

%%%%-----导入宏包-----%%%%
\usepackage{style/shubeamer}   % SHU 主题（配色、帧标题、页脚、Python代码样式）
\usepackage{amsmath,amsfonts,amssymb,bm}
\usepackage{graphicx,hyperref,url}
\usepackage{subcaption}
\usepackage{booktabs}
\usepackage[super,square]{natbib}
\usepackage{wrapfig}          % 文字绕排
\usepackage{calligra}         % 结尾页花体字

%%%%-----设置字体-----%%%%
\setmainfont{CMU Serif}
\setsansfont{CMU Sans Serif}

%%%%-----常用自定义颜色（正文可直接使用）-----%%%%
\definecolor{nvidia}{RGB}{102,156,28}
\definecolor{red}{RGB}{184,13,73}
\definecolor{orange}{RGB}{242,151,36}
\definecolor{darkteal}{RGB}{43,106,108}
\definecolor{darkblue}{RGB}{4,37,58}

%%%%-----每节开始时自动显示当前节目录-----%%%%
\AtBeginSection[]
{
	\begin{frame}<beamer>
		\frametitle{\textbf{目录}}
		\tableofcontents[currentsection,sectionstyle=show/hide,subsectionstyle=hide]
	\end{frame}
}

%%%%-----设定图片存放路径-----%%%%
\graphicspath{{figures/}}

%%%%-----首页信息设置-----%%%%
\AtBeginDocument{%
	%%%%----短标题[显示在页脚]{封面标题}
	\title[Python金融数据分析 · 第4章]{\fontsize{16pt}{22pt}\selectfont {Python金融数据分析}}
	%%%%----副标题（本章主题）
	\subtitle{\fontsize{9pt}{14pt}\selectfont \textit{第4章\quad pandas数据分析}}
	%%%%----个人信息
	\author[王伟冠, 吕文俊]{%
		\begin{tabular}{ll}%
			主讲人： & 王伟冠, 吕文俊\tabularnewline%
			院\quad 系： & 经济学院金融系\tabularnewline
		\end{tabular}
	}
	%%%%----机构信息
	\institute[SHU]{上海大学}
	%%%%----日期信息
	\date[\today]{\today}}

%%%%-----数学宏（向量、矩阵、算子等，见 style/macros.tex）-----%%%%
\input{style/macros.tex}


\begin{document}

%% 封面
\begin{frame}
	\titlepage
\end{frame}

%% 目录页
\section*{目录}
\begin{frame}
	\frametitle{\textbf{目录}}
	\begin{columns}[T]\column{.49\textwidth}\tableofcontents[sections={2-7}]\column{.49\textwidth}\tableofcontents[sections={8-13}]\end{columns}
\end{frame}


%==============================================================
\section{初识pandas}
%==============================================================

% 多周教学路线
\begin{frame}{教学路线：可分四个单元推进}
\begin{enumerate}
\item 表与索引：选择、赋值、对齐，先说明一行代表什么。
\item 清洗：类型、重复、缺失、异常，保留处理依据。
\item 合并与变形：连接关系、行数变化、长表与宽表。
\item 分组与应用：聚合、组内计算、订单日报与金融迁移。
\end{enumerate}
\begin{block}{课堂节奏}
每单元可安排一周，也可将讲解、上机和讲评分开。完成本单元练习后再推进，不要求一周讲完整章。
\end{block}
\end{frame}
\begin{frame}{配套 Notebook 与学习顺序}
\begin{itemize}
\item 单元一：\texttt{04\_pandas} 第1--5节。
\item 单元二：\texttt{04\_2\_Cleaning}，整理成绩登记。
\item 单元三：\texttt{04\_3\_Wrangling}，合并订单与商品。
\item 单元四：\texttt{04\_4\_GroupBy}，完成订单日报；\\
再回到 \texttt{04\_pandas} 第6--9节，完成金融案例。
\end{itemize}
\begin{block}{运行方式}
三个补充 Notebook 均可单独从头运行。返回主 Notebook 时从头执行，恢复本文件的变量。
\end{block}
\end{frame}

\begin{frame}{本章学习目标}
	\begin{itemize}
		\item 理解 Series 与 DataFrame 两种核心结构
		\item 掌握数据的选取：按列、按标签 \texttt{loc}、按位置 \texttt{iloc}、条件筛选
		\item 掌握排序与分组聚合（\texttt{groupby}）
		\item 掌握金融时间序列的基本操作：日期索引、\texttt{pct\_change}、重采样
		\item 会处理缺失值，会用 \texttt{read\_csv} 读入真实数据
		\item 掌握清洗、连接、长宽表转换与组内统计，能够对一张股价表完成"收益率—风险—相关性"的完整分析
	\end{itemize}
	\begin{block}{前置知识}
		第2章的字典/列表、第3章的NumPy——pandas 就建在 NumPy 之上。
	\end{block}
\end{frame}

\begin{frame}{NumPy不够用了？}
	\begin{itemize}
		\item 第3章我们算出了收益率数组，但马上会遇到新问题：
		\begin{itemize}
			\item 每个数字对应\alert{哪一天}？——NumPy 数组没有标签；
			\item 行情数据里有日期、代码、名称——普通数组通常使用统一的数据类型；
			\item 数据有\alert{缺失}（停牌）怎么办？
			\item 从 Excel/CSV 读进来的表格如何直接分析？
		\end{itemize}
		\item \textbf{pandas} 就是为此而生：给数组加上行列标签，成为"数据表"。
	\end{itemize}
	\begin{exampleblock}{一个直观比喻}
		NumPy 数组像一张写满数字的草稿纸；\\
		pandas DataFrame 像一张\alert{带表头和日期列的 Excel 表}，而且能用代码操作。
	\end{exampleblock}
\end{frame}

\begin{frame}[fragile]{导入约定与两种核心结构}
	\begin{lstlisting}
import numpy as np
import pandas as pd        # 约定俗成，照抄
\end{lstlisting}
	\begin{table}
		\centering
		\begin{tabular}{lll}
			\toprule
			结构 & 可以理解为 & 金融中的例子 \\
			\midrule
			\texttt{Series} & 一列带索引的数据 & 某股票的收盘价序列 \\
			\texttt{DataFrame} & 一张带行列标签的表 & 多股票$\times$多日期 的价格表 \\
			\bottomrule
		\end{tabular}
	\end{table}
	\begin{itemize}
		\item DataFrame 本质上就是"共享同一个索引的若干 Series"。
		\item 本章几乎所有内容都围绕这张"表"的查看、选取、变形与统计。
	\end{itemize}
\end{frame}


%==============================================================
\section{Series：带标签的一列数据}
%==============================================================

\begin{frame}[fragile]{创建Series}
	\begin{lstlisting}
s = pd.Series([1500.35, 128.60, 55.20],
              index=["600519", "000858", "601318"])
s
\end{lstlisting}
	\begin{lstlisting}
600519    1500.35
000858     128.60
601318      55.20
dtype: float64
\end{lstlisting}
	\begin{itemize}
		\item 左边是\textbf{索引}（index，相当于标签），右边是值。
		\item 不指定索引时，默认为 0, 1, 2, \ldots（像列表）。
		\item 索引可以是股票代码、日期、任何你需要的标签。
	\end{itemize}
\end{frame}

\begin{frame}[fragile]{按标签取值与向量化}
	\begin{lstlisting}
s["600519"]            # 1500.35  按标签取值
s[["600519", "601318"]]  # 取多个

s * 1.05               # 所有值涨5%（向量化，继承自NumPy）
s[s > 100]             # 筛选出大于100的
\end{lstlisting}
	\begin{itemize}
		\item Series 支持 NumPy 式向量化运算与布尔筛选——旧知识直接复用。
		\item 取值时用的是\alert{标签}而不是位置：\texttt{s["600519"]}。
	\end{itemize}
	\begin{alertblock}{关键概念：对齐}
		两个 Series 运算时按\textbf{索引对齐}，而不是按位置：\\
		"茅台的价格"只和"茅台的权重"相乘，顺序打乱也不会错。
	\end{alertblock}
\end{frame}


%==============================================================
\section{DataFrame：金融数据的"主战场"}
%==============================================================

\begin{frame}[fragile]{从字典创建DataFrame}
	\begin{lstlisting}
data = {"茅台": [1500.0, 1520.0, 1515.0],
        "五粮液": [128.0, 129.5, 127.0],
        "平安": [55.0, 54.5, 56.0]}
df = pd.DataFrame(data,
                  index=["2026-07-01", "2026-07-02", "2026-07-03"])
df
\end{lstlisting}
	\begin{itemize}
		\item 字典的\textbf{键}变成列名，列表里的值成为该列数据。
		\item 约定俗成的金融数据方向：\alert{行=日期，列=资产}。
		\item Notebook 中直接写变量名 \texttt{df}，会以漂亮的表格显示。
	\end{itemize}
\end{frame}

\begin{frame}[fragile]{查看数据的基本动作}
	\begin{lstlisting}
df.shape       # (3, 3)      行数、列数
df.columns     # 列名：Index(['茅台', '五粮液', '平安'])
df.index       # 行索引（这里是日期）
df.head(2)     # 前2行（head()默认5行）
df.tail(1)     # 最后1行
df.describe()  # 每列的统计摘要：均值、标准差、分位数等
\end{lstlisting}
	\begin{itemize}
		\item 拿到任何新数据，先把这些"体检动作"做一遍。
		\item \texttt{describe()} 相当于免费的初步报告，务必常用。
	\end{itemize}
\end{frame}


%==============================================================
% 自编教学补充；阅读来源见本章 README。
\begin{frame}[fragile]{3A 先说明一行代表什么}
\small 这张表每行代表一次学生成绩记录，学号是标识而非数量。金融表也可以每行代表“股票—日期”，不必总是每列一只股票。
\medskip
\begin{lstlisting}
students = pd.DataFrame({
    "学号": ["001", "002", "003", "004"],
    "姓名": ["小林", "小周", "小陈", "小许"],
    "班级": ["A", "A", "B", "B"],
    "成绩": [82, 95, 68, 88]})
students.info()
students.head()
\end{lstlisting}
\end{frame}

\begin{frame}[fragile]{3B 一列与一列表：维度不同}
\small 单层方括号取一列得到 Series；列名列表得到 DataFrame。先观察 type 和 shape，再决定后面如何操作。
\medskip
\begin{lstlisting}
one_column = students["成绩"]
one_table = students[["成绩"]]
print(type(one_column), one_column.shape)
print(type(one_table), one_table.shape)
print(students.describe())
\end{lstlisting}
\end{frame}

\section{数据选取与筛选}
%==============================================================

\begin{frame}[fragile]{取列与取行}
	\begin{lstlisting}
df["茅台"]              # 取一列 -> Series
df[["茅台", "平安"]]    # 取多列 -> DataFrame
df["2026-07-01"]        # 报错！方括号默认是"按列"
\end{lstlisting}
	\begin{itemize}
		\item \texttt{df[...]} 里放列名：取\textbf{列}。
		\item 想按行取，用切片（按位置）或下一页的 \texttt{loc/iloc}。
	\end{itemize}
	\begin{lstlisting}
df[0:2]          # 前两行（位置切片，含头不含尾）
\end{lstlisting}
	\begin{alertblock}{最常见的困惑}
		为什么 \texttt{df["茅台"]} 可以，\texttt{df["2026-07-01"]} 却报错？\\
		因为 \texttt{[]} 默认针对\textbf{列}。规则记牢，少查一半报错。
	\end{alertblock}
\end{frame}

\begin{frame}[fragile]{loc：按标签选；iloc：按位置选}
	\begin{lstlisting}
df.loc["2026-07-02"]              # 该行全部列
df.loc["2026-07-02", "茅台"]      # 指定行、指定列 -> 1520.0
df.loc[:, "茅台"]                 # 所有行、茅台列

df.iloc[0]                        # 第0行（按位置）
df.iloc[0:2, 1]                   # 前2行、第1列
df.iloc[-1]                       # 最后一行
\end{lstlisting}
	\begin{itemize}
		\item \texttt{loc}：用\textbf{标签}（日期、列名）；写法 \texttt{[行, 列]}。
		\item \texttt{iloc}：用\textbf{位置}（0, 1, 2\ldots）；与列表索引规则相同。
		\item \alert{先想清楚"我要按名字找还是按位置找"，再选工具}。
	\end{itemize}
\end{frame}

\begin{frame}[fragile]{条件筛选：选出满足条件的行}
	\begin{lstlisting}
df[df["茅台"] > 1510]       # 茅台价格高于1510的日子
df[(df["茅台"] > 1500) & (df["平安"] > 55)]
\end{lstlisting}
	\begin{itemize}
		\item 与 NumPy 布尔索引完全一致：先算出布尔列，再用来筛行。
		\item 多条件组合用 \texttt{\&}、\texttt{|}，每个条件加括号。
		\item 金融例子：涨幅超过5\%的日子、PE低于10且股息率高于4\%的股票。
	\end{itemize}
	\begin{exampleblock}{课堂练习}
		用一行代码选出"五粮液价格低于128"的所有交易日。
	\end{exampleblock}
\end{frame}


%==============================================================
% 自编教学补充；阅读来源见本章 README。
\begin{frame}[fragile]{4A 标签切片与位置切片}
\small 学号设为索引后，loc 用学号找人，iloc 按当前位置找人。在此唯一且有序的索引上，loc 切片包括末端标签，iloc 不包括末端位置。
\medskip
\begin{lstlisting}
by_id = students.set_index("学号")
print(by_id.loc["001":"003", ["姓名", "成绩"]])
print(by_id.iloc[0:2, [0, 2]])
print(by_id.reset_index().columns.tolist())
\end{lstlisting}
\end{frame}

\begin{frame}[fragile]{4B 条件筛选与安全赋值}
\small 用 loc 一次指定行与列进行赋值；不要在 df[条件][列名] 的临时结果上修改原表。保留原始成绩，另建调整列。
\medskip
\begin{lstlisting}
students["调整成绩"] = students["成绩"]
mask = students["班级"].eq("B")
students.loc[mask, "调整成绩"] = students.loc[mask, "成绩"] + 2
selected = students.loc[students["成绩"].ge(80),
                        ["姓名", "班级", "成绩"]].copy()
print(selected)
\end{lstlisting}
\end{frame}
\begin{frame}[fragile]{课堂练习：4B 条件筛选与安全赋值}
筛选 A 班且成绩至少为 90 的学生，仅展示姓名与成绩。
\medskip
\begin{block}{检查思路}
先写出预期结果，再运行。参考答案见配套 Notebook 的折叠单元格。
\end{block}
\end{frame}

\begin{frame}[fragile]{4C 对齐：标签相同才相加}
\small Series 运算按标签匹配。缺少对应项会得到缺失值；fill\_value=0 只适合“未出现确实表示没有”的业务含义。
\medskip
\begin{lstlisting}
morning = pd.Series([3, 5], index=["茶", "咖啡"])
afternoon = pd.Series([2, 4], index=["咖啡", "果汁"])
print(morning + afternoon)
print(morning.add(afternoon, fill_value=0))
\end{lstlisting}
\end{frame}

\section{排序与分组}
%==============================================================

\begin{frame}[fragile]{排序}
	\begin{lstlisting}
df.sort_values("茅台", ascending=False)  # 按茅台价格降序
df.sort_index()                          # 按行索引(日期)排序
\end{lstlisting}
	\begin{itemize}
		\item \texttt{sort\_values(列名)}：按某列的值给\textbf{行}排序。
		\item \texttt{sort\_index()}：按行索引排序——日期数据常用。
		\item 排序返回\textbf{新表}，原表不变（除非加 \texttt{inplace=True}）。
	\end{itemize}
	\begin{alertblock}{习惯}
		时间序列数据到手先 \texttt{sort\_index()}，保证日期升序，后续计算才安全。
	\end{alertblock}
\end{frame}

\begin{frame}[fragile]{groupby：分组聚合}
	\begin{lstlisting}
stocks = pd.DataFrame({
    "代码":  ["600519", "000858", "601318", "600036"],
    "行业":  ["白酒", "白酒", "保险", "银行"],
    "市盈率": [28.5, 18.2, 8.1, 6.3]})

stocks.groupby("行业")["市盈率"].mean()
\end{lstlisting}
	\begin{lstlisting}
行业
保险     8.10
银行     6.30
白酒    23.35
\end{lstlisting}
	\begin{itemize}
		\item 三步走：\textbf{分组（groupby）$\to$ 选列 $\to$ 聚合（mean/sum/max\ldots）}。
		\item 回答"每个行业的平均市盈率"这类问题，一行搞定。
	\end{itemize}
\end{frame}


%==============================================================
\section{数据清洗与输入输出}
% 自编教学补充；阅读来源见本章 README。
\begin{frame}[fragile]{1. 原始登记表：保留文本才能看见问题}
\small 教学合成数据，每行一次成绩登记。故意保留首尾空格、重复记录、缺考和不合理成绩；学号按字符串读入以保留前导零。
\medskip
\begin{lstlisting}
from io import StringIO
import pandas as pd
raw_csv = """学号,姓名,班级,成绩,日期
001, 小林 ,A,82,2026-09-01
002,小周,a,95,2026-09-01
003,小陈,B,缺考,2026-09-02
004,小许,B,108,2026-09-02
001, 小林 ,A,82,2026-09-01
005,小吴, B ,76,日期待核"""
raw = pd.read_csv(StringIO(raw_csv), dtype="string")
raw
\end{lstlisting}
\end{frame}

\begin{frame}[fragile]{2. 先检查，再清洗}
\small info 看类型和非缺失数，duplicated 找整行重复。这里“缺考”是普通文本，尚不会被 isna 自动识别。
\medskip
\begin{lstlisting}
raw.info()
print(raw.isna().sum())
print("整行重复数：", raw.duplicated().sum())
raw.loc[raw.duplicated(keep=False)]
\end{lstlisting}
\end{frame}

\begin{frame}[fragile]{3. str 方法：把单个字符串操作扩展到整列}
\small copy 保留原始登记。统一班级格式后再分组，否则 A 和 a 会被算成两个班。
\medskip
\begin{lstlisting}
clean = raw.copy()
clean["姓名"] = clean["姓名"].str.strip()
clean["班级"] = clean["班级"].str.strip().str.upper()
print(clean["班级"].value_counts())
print(clean["姓名"].str.contains("小", na=False))
\end{lstlisting}
\end{frame}

\begin{frame}[fragile]{4. 类型转换：转换失败的行必须可追踪}
\small 先把明确的缺考编码改成缺失，再转数值。coerce 将无法解析的值变成缺失；保留原列才能区分缺考与格式错误。
\medskip
\begin{lstlisting}
score_text = clean["成绩"].replace("缺考", pd.NA)
clean["分数"] = pd.to_numeric(score_text, errors="coerce")
clean["登记日"] = pd.to_datetime(clean["日期"], errors="coerce")
parse_failed = score_text.notna() & clean["分数"].isna()
print(clean.loc[parse_failed | clean["登记日"].isna()])
print(clean.dtypes)
\end{lstlisting}
\end{frame}

\begin{frame}[fragile]{5. 重复键不等于重复行}
\small 本例每名学生只有一次考试，学号应唯一。已确认整行相同的重复登记可以删除；若同一学号分数不同，应先核实，不能任意保留。
\medskip
\begin{lstlisting}
print(clean.loc[clean.duplicated("学号", keep=False)])
clean = clean.drop_duplicates().reset_index(drop=True)
print("清理后记录数：", len(clean))
print("学号唯一：", clean["学号"].is_unique)
\end{lstlisting}
\end{frame}

\begin{frame}[fragile]{6. 异常值与缺失值分开记录}
\small 考试满分 100，因此 108 是待核异常，不能直接截成 100。缺考也不等于 0 分。设置标记后，在统计列中排除异常。
\medskip
\begin{lstlisting}
clean["成绩异常"] = (clean["分数"].notna() &
                       ~clean["分数"].between(0, 100))
clean["有效成绩"] = clean["分数"].mask(clean["成绩异常"])
print(clean[["学号", "成绩", "成绩异常", "有效成绩"]])
print("有效成绩均值：", clean["有效成绩"].mean())
\end{lstlisting}
\end{frame}
\begin{frame}[fragile]{课堂练习：6. 异常值与缺失值分开记录}
计算有效成绩数量、缺考数量、异常成绩数量。为什么不能把三者混成一类？
\medskip
\begin{block}{检查思路}
先写出预期结果，再运行。参考答案见配套 Notebook 的折叠单元格。
\end{block}
\end{frame}

\begin{frame}[fragile]{7. 删除与填充：输出改变了什么}
\small 比较同一列的有效值平均与填零平均。dropna 的 subset 明确本次分析依赖哪些列；不能因为日期缺失就删除不依赖日期的成绩。
\medskip
\begin{lstlisting}
print(clean["有效成绩"].mean())
print(clean["有效成绩"].fillna(0).mean())
usable_scores = clean.dropna(subset=["有效成绩"])
print(usable_scores[["学号", "有效成绩"]])
\end{lstlisting}
\end{frame}

\begin{frame}[fragile]{8. 分类与映射：明确区间边界}
\small cut 按固定阈值分段，right=False 表示左闭右开。100 包含在 [90,101) 内；没有有效成绩的行仍保留缺失等级。
\medskip
\begin{lstlisting}
clean["等级"] = pd.cut(clean["有效成绩"],
    bins=[0, 60, 80, 90, 101], right=False,
    labels=["待提高", "合格", "良好", "优秀"])
clean["班级名称"] = clean["班级"].map({"A": "一班", "B": "二班"})
clean[["学号", "有效成绩", "等级", "班级名称"]]
\end{lstlisting}
\end{frame}

\begin{frame}[fragile]{9. 导出与读回：CSV 不保存所有类型}
\small 先在内存中演示 CSV 往返。读回时显式指定学号类型；分类类型和日期类型需要重建。真实文件路径可替换 StringIO。
\medskip
\begin{lstlisting}
csv_text = clean.to_csv(index=False)
restored = pd.read_csv(StringIO(csv_text), dtype={"学号": "string"})
print(restored["学号"].tolist())
print(len(raw), len(clean), len(restored))
assert len(raw) == 6 and len(clean) == 5
assert restored["学号"].iloc[0] == "001"
\end{lstlisting}
\end{frame}

\begin{frame}[fragile]{10. 本单元练习：给出可解释的清洗记录}
\small 交付三项结果：保留几行、排除哪些成绩、哪些日期待核。下一单元学习把这类记录与另一张信息表合并。
\medskip
\begin{lstlisting}
audit = pd.Series({
    "输入行数": len(raw), "保留行数": len(clean),
    "有效成绩数": clean["有效成绩"].count(),
    "异常成绩数": clean["成绩异常"].sum(),
    "待核日期数": clean["登记日"].isna().sum()})
print(audit)
\end{lstlisting}
\end{frame}
\begin{frame}[fragile]{课堂练习：10. 本单元练习：给出可解释的清洗记录}
找出有效成绩低于 80 的学生，保留学号、姓名、班级、有效成绩。解释为什么没有把缺考者选进来。
\medskip
\begin{block}{检查思路}
先写出预期结果，再运行。参考答案见配套 Notebook 的折叠单元格。
\end{block}
\end{frame}

\section{合并与长宽表转换}
% 自编教学补充；阅读来源见本章 README。
\begin{frame}[fragile]{1. 两张表的粒度与键}
\small 合成订单表每行一笔订单，商品表每行一种商品。同一商品可以出现在多笔订单中，因此合并关系是多对一。
\medskip
\begin{lstlisting}
import pandas as pd
orders = pd.DataFrame({
    "订单": [1, 2, 3, 4], "商品": ["T", "C", "T", "X"],
    "数量": [2, 1, 3, 1]})
products = pd.DataFrame({
    "商品": ["T", "C", "J"],
    "名称": ["茶", "咖啡", "果汁"], "单价": [10, 20, 15]})
print(orders)
print(products)
\end{lstlisting}
\end{frame}

\begin{frame}[fragile]{2. 左连接：保留每笔订单}
\small on 明确匹配键；left 保留左表订单。indicator 标记未匹配行，validate 检查商品表的键唯一，避免合并后订单被复制。
\medskip
\begin{lstlisting}
joined = orders.merge(products, on="商品", how="left",
    validate="many_to_one", indicator=True)
print(joined)
print(joined["_merge"].value_counts())
assert len(joined) == len(orders)
\end{lstlisting}
\end{frame}

\begin{frame}[fragile]{3. 内连接与外连接：观察谁消失、谁新增}
\small inner 只保留双方存在的商品匹配；outer 还会显示没有订单的商品。未匹配不代表金额为零，应单独列出。
\medskip
\begin{lstlisting}
inner = orders.merge(products, on="商品", how="inner")
outer = orders.merge(products, on="商品", how="outer",
                     indicator=True)
print(len(inner), len(outer))  # 3, 5
print(outer[["订单", "商品", "_merge"]])
\end{lstlisting}
\end{frame}
\begin{frame}[fragile]{课堂练习：3. 内连接与外连接：观察谁消失、谁新增}
哪笔订单在 inner 中消失？哪个商品只在右表中？如果要核对全部订单，该选哪种连接？
\medskip
\begin{block}{检查思路}
先写出预期结果，再运行。参考答案见配套 Notebook 的折叠单元格。
\end{block}
\end{frame}

\begin{frame}[fragile]{4. 重复键：合并可能让行数变多}
\small 在商品表中故意重复 T，观察两笔 T 订单各匹配两次。解决方法是核实键与版本，而不是合并后随意去重。
\medskip
\begin{lstlisting}
bad_products = pd.concat([products, products.iloc[[0]]],
                          ignore_index=True)
expanded = orders.merge(bad_products, on="商品", how="left")
print(len(orders), len(expanded))  # 4, 6
print(bad_products.duplicated("商品").sum())
# 加 validate="many_to_one" 将直接报错，课堂可自行尝试。
\end{lstlisting}
\end{frame}

\begin{frame}[fragile]{5. concat：把同结构记录上下接起来}
\small concat 不按商品键寻找匹配。ignore\_index 重建行号，不检查业务编号是否重复；拼接后仍要检查订单号。
\medskip
\begin{lstlisting}
next_orders = pd.DataFrame({
    "订单": [5, 6], "商品": ["C", "T"], "数量": [2, 1]})
all_orders = pd.concat([orders, next_orders], ignore_index=True)
print(all_orders)
print(all_orders["订单"].is_unique)
\end{lstlisting}
\end{frame}

\begin{frame}[fragile]{6. concat 的另一方向：横向按索引对齐}
\small axis=1 横向拼接时按索引标签对齐，不会因为行数相同就按位置贴在一起。先检查索引含义。
\medskip
\begin{lstlisting}
left = pd.Series([80, 90], index=["001", "002"], name="数学")
right = pd.Series([70, 85], index=["002", "001"], name="英语")
print(pd.concat([left, right], axis=1))
\end{lstlisting}
\end{frame}

\begin{frame}[fragile]{7. 长表转宽表：pivot 不做统计}
\small 长表每行是一名学生的一门课成绩；宽表每行一名学生、每列一门课。pivot 要求“学号—课程”组合唯一。
\medskip
\begin{lstlisting}
long_scores = pd.DataFrame({
    "学号": ["001", "001", "002", "002"],
    "课程": ["数学", "英语", "数学", "英语"],
    "成绩": [80, 90, 70, 85]})
wide = long_scores.pivot(index="学号", columns="课程", values="成绩")
print(wide)
\end{lstlisting}
\end{frame}

\begin{frame}[fragile]{8. 宽表转长表：melt 保留身份列}
\small id\_vars 指定不被展开的标识列。melt 将列名变成课程字段，将原来的单元格变成成绩字段。
\medskip
\begin{lstlisting}
back = wide.reset_index().melt(id_vars="学号",
    var_name="课程", value_name="成绩")
print(back.sort_values(["学号", "课程"]))
assert len(back) == len(long_scores)
\end{lstlisting}
\end{frame}

\begin{frame}[fragile]{9. 重复组合：先决定是否需要聚合}
\small 加入一次补考后，同一学生同一课程出现两次。pivot\_table 可以聚合；取最大值代表“保留最高成绩”，这是业务规则。
\medskip
\begin{lstlisting}
retake = pd.DataFrame({"学号": ["002"],
                       "课程": ["数学"], "成绩": [88]})
attempts = pd.concat([long_scores, retake], ignore_index=True)
best = attempts.pivot_table(index="学号", columns="课程",
    values="成绩", aggfunc="max")
print(best)
\end{lstlisting}
\end{frame}
\begin{frame}[fragile]{课堂练习：9. 重复组合：先决定是否需要聚合}
将 max 改为 mean，002 的数学成绩变成多少？两种答案分别回答什么问题？
\medskip
\begin{block}{检查思路}
先写出预期结果，再运行。参考答案见配套 Notebook 的折叠单元格。
\end{block}
\end{frame}

\begin{frame}[fragile]{10. 多层索引：把组合键显式写出来}
\small MultiIndex 可以表示“学号—课程”这类组合标签。先学 set\_index 和 reset\_index，再选学 stack、unstack。
\medskip
\begin{lstlisting}
indexed = long_scores.set_index(["学号", "课程"]).sort_index()
print(indexed.loc[("001", "数学"), "成绩"])
print(indexed.unstack("课程"))
print(indexed.reset_index())
\end{lstlisting}
\end{frame}

\begin{frame}[fragile]{11. 综合检查：只汇总匹配成功的金额}
\small 价格未知的 X 订单不能当作零金额。先列出未匹配订单，再汇总可计算的部分，并说明覆盖范围。
\medskip
\begin{lstlisting}
matched = joined.loc[joined["_merge"] == "both"].copy()
matched["金额"] = matched["数量"] * matched["单价"]
print(matched[["订单", "名称", "金额"]])
print("已知金额合计：", matched["金额"].sum())  # 70
print("未匹配订单数：", (joined["_merge"] == "left_only").sum())
\end{lstlisting}
\end{frame}
\begin{frame}[fragile]{课堂练习：11. 综合检查：只汇总匹配成功的金额}
将商品 X 补录为“水”、单价 5，重新合并全部原订单。总金额应是多少？
\medskip
\begin{block}{检查思路}
先写出预期结果，再运行。参考答案见配套 Notebook 的折叠单元格。
\end{block}
\end{frame}

\section{分组统计与综合应用}
% 自编教学补充；阅读来源见本章 README。
\begin{frame}[fragile]{1. 校园饮品订单：一行一笔订单}
\small 六笔合成订单用于手工核对。金额为数量乘单价；本例价格均已匹配，日期是实际有订单的日期。
\medskip
\begin{lstlisting}
import pandas as pd
sales = pd.DataFrame({
    "订单": [1, 2, 3, 4, 5, 6],
    "日期": ["2026-09-01"] * 3 + ["2026-09-02"] * 3,
    "门店": ["东", "东", "西", "东", "西", "西"],
    "商品": ["茶", "咖啡", "茶", "茶", "咖啡", "茶"],
    "数量": [2, 1, 1, 3, 2, 2], "单价": [10, 20, 10, 10, 20, 10]})
sales["日期"] = pd.to_datetime(sales["日期"])
sales["金额"] = sales["数量"] * sales["单价"]
sales
\end{lstlisting}
\end{frame}

\begin{frame}[fragile]{2. groupby：拆分、计算、合并}
\small 先按门店拆成小表，在各小表内求和，再组成汇总表。as\_index=False 使门店保留为普通列，便于继续合并。
\medskip
\begin{lstlisting}
store_total = sales.groupby("门店", as_index=False)["金额"].sum()
print(store_total)
print(sales.groupby(["门店", "商品"])["金额"].sum())
\end{lstlisting}
\end{frame}

\begin{frame}[fragile]{3. 命名聚合：明确每个指标的含义}
\small 总金额、订单数和平均每单金额是不同指标。命名聚合直接给结果列起名；金额均值的分母是订单行数。
\medskip
\begin{lstlisting}
summary = sales.groupby("门店", as_index=False).agg(
    总金额=("金额", "sum"),
    订单数=("订单", "size"),
    总杯数=("数量", "sum"),
    平均每单金额=("金额", "mean"))
print(summary)
\end{lstlisting}
\end{frame}

\begin{frame}[fragile]{4. size 与 count：分母不能混淆}
\small size 统计组内行数，count 统计某列非缺失值数。先明确要数订单、有效金额还是不同顾客，再选择统计方法。
\medskip
\begin{lstlisting}
with_missing = sales.copy()
with_missing.loc[0, "金额"] = float("nan")
counts = with_missing.groupby("门店").agg(
    行数=("金额", "size"), 有效金额数=("金额", "count"))
print(counts)
\end{lstlisting}
\end{frame}
\begin{frame}[fragile]{课堂练习：4. size 与 count：分母不能混淆}
东店的行数与有效金额数各是多少？金额缺失能否当作免费订单？
\medskip
\begin{block}{检查思路}
先写出预期结果，再运行。参考答案见配套 Notebook 的折叠单元格。
\end{block}
\end{frame}

\begin{frame}[fragile]{5. 加权均值：平均每杯价格}
\small 平均每单金额不等于平均每杯价格。后者用总金额除以总杯数，相当于按购买数量给单价加权。
\medskip
\begin{lstlisting}
summary["平均每杯价格"] = summary["总金额"] / summary["总杯数"]
print(summary[["门店", "平均每单金额", "平均每杯价格"]])
print("全体每杯均价：", sales["金额"].sum() / sales["数量"].sum())
\end{lstlisting}
\end{frame}

\begin{frame}[fragile]{6. transform：把组内统计带回原行}
\small agg 每组输出一行，transform 返回与原表等长且对齐的结果。用每笔金额除以所属门店总金额得到组内占比。
\medskip
\begin{lstlisting}
sales["门店总金额"] = sales.groupby("门店")["金额"].transform("sum")
sales["门店内占比"] = sales["金额"] / sales["门店总金额"]
print(sales[["订单", "门店", "金额", "门店内占比"]])
print(sales.groupby("门店")["门店内占比"].sum())
\end{lstlisting}
\end{frame}

\begin{frame}[fragile]{7. 排名与组内前两名}
\small 先按金额降序，再按订单号升序决定并列次序，组内取前两行。结果是两笔订单，而不一定是两个不同金额档位。
\medskip
\begin{lstlisting}
ordered = sales.sort_values(["门店", "金额", "订单"],
                             ascending=[True, False, True])
top2 = ordered.groupby("门店").head(2)
print(top2[["订单", "门店", "金额"]])
sales["组内名次"] = sales.groupby("门店")["金额"].rank(
    ascending=False, method="min")
\end{lstlisting}
\end{frame}

\begin{frame}[fragile]{8. 交叉表：计数与金额表不同}
\small crosstab 默认数记录。pivot\_table 显式指定金额与 sum 才是金额汇总。本例已知交易明细完整，因此没有订单的组合可填零。
\medskip
\begin{lstlisting}
print(pd.crosstab(sales["门店"], sales["商品"]))
amount_table = sales.pivot_table(index="门店", columns="商品",
    values="金额", aggfunc="sum", fill_value=0)
print(amount_table)
\end{lstlisting}
\end{frame}

\begin{frame}[fragile]{9. 按日汇总与滚动窗口}
\small 先汇总到“日”，再滚动才是相邻观测日的窗口。rolling(2) 是两行，不自动补齐日历中没有记录的日期。
\medskip
\begin{lstlisting}
daily = sales.groupby("日期")["金额"].sum().sort_index()
print(daily)
print(daily.rolling(2, min_periods=2).mean())
print(sales.set_index("日期")["金额"].resample("D").sum(min_count=1))
\end{lstlisting}
\end{frame}

\begin{frame}[fragile]{10. 金融迁移：按股票计算收益率}
\small 合成长表每行一只股票一天的收盘价。先按股票和日期排序，在每只股票内部错位，避免把另一只股票的价格当作昨日价。
\medskip
\begin{lstlisting}
quotes = pd.DataFrame({
    "股票": ["A", "B", "A", "B", "A", "B"],
    "日期": pd.to_datetime(["2026-09-01"] * 2 +
                          ["2026-09-02"] * 2 + ["2026-09-03"] * 2),
    "收盘价": [100, 50, 102, 49, 101, 50]})
quotes = quotes.sort_values(["股票", "日期"])
quotes["昨日价"] = quotes.groupby("股票")["收盘价"].shift(1)
quotes["收益率"] = quotes["收盘价"] / quotes["昨日价"] - 1
print(quotes)
\end{lstlisting}
\end{frame}

\begin{frame}[fragile]{11. 金融迁移：从长表到收益率矩阵}
\small pivot 后每列一只股票；pct\_change(fill\_method=None) 不自动填补缺失价格。比较两种写法的同一观测值。
\medskip
\begin{lstlisting}
price_matrix = quotes.pivot(index="日期", columns="股票", values="收盘价")
return_matrix = price_matrix.pct_change(fill_method=None)
print(return_matrix)
expected = quotes.loc[(quotes["股票"] == "A") &
    (quotes["日期"] == "2026-09-02"), "收益率"].iloc[0]
assert abs(return_matrix.loc["2026-09-02", "A"] - expected) < 1e-12
\end{lstlisting}
\end{frame}

\begin{frame}[fragile]{12. 综合练习与验收}
\small 先用订单手算，再检查汇总表。全表金额应为 140 元、总杯数为 11；两门店各 70 元，但每杯均价不同。
\medskip
\begin{lstlisting}
assert sales["金额"].sum() == 140
assert sales["数量"].sum() == 11
assert summary["总金额"].sum() == sales["金额"].sum()
print(summary)
\end{lstlisting}
\end{frame}
\begin{frame}[fragile]{课堂练习：12. 综合练习与验收}
按“日期—门店”统计金额与杯数，并计算每杯均价。哪一天全体销售金额更高？
\medskip
\begin{block}{检查思路}
先写出预期结果，再运行。参考答案见配套 Notebook 的折叠单元格。
\end{block}
\end{frame}

\begin{frame}[fragile]{13. 综合案例起点：从原始订单到日报}
\small 下面重新开始一套独立的合成数据。每行原本应是一笔订单，但有重复登记、数量缺失和未知商品；目标是输出可核对的已知金额日报。
\medskip
\begin{lstlisting}
raw_orders = pd.DataFrame({
    "订单": ["01", "02", "03", "04", "05", "06", "01"],
    "日期": ["2026-09-01"] * 3 + ["2026-09-02"] * 3 + ["2026-09-01"],
    "门店": ["东", "东", "西", "东", "西", "西", "东"],
    "商品": [" t ", "C", "T", "T", "C", "X", " t "],
    "数量": ["2", "1", "缺失", "3", "2", "1", "2"]})
catalog = pd.DataFrame({"商品": ["T", "C"], "单价": [10, 20]})
raw_orders
\end{lstlisting}
\end{frame}

\begin{frame}[fragile]{14. 清洗：原始值与统计值并存}
\small 保留 raw\_orders。先统一商品代码，再删除完全相同的登记，转换日期和数量；原始数量列保留，便于回查缺失原因。
\medskip
\begin{lstlisting}
tidy_orders = raw_orders.copy()
tidy_orders["商品"] = tidy_orders["商品"].str.strip().str.upper()
tidy_orders = tidy_orders.drop_duplicates().reset_index(drop=True)
tidy_orders["日期"] = pd.to_datetime(tidy_orders["日期"])
tidy_orders["有效数量"] = pd.to_numeric(tidy_orders["数量"], errors="coerce")
assert tidy_orders["订单"].is_unique
print(tidy_orders)
\end{lstlisting}
\end{frame}

\begin{frame}[fragile]{15. 合并：把待核订单单独列出}
\small 左连接保留每笔订单。已匹配商品且数量为正才进入本例金额统计；数量缺失和价格未知都保留在待核清单中。
\medskip
\begin{lstlisting}
enriched = tidy_orders.merge(catalog, on="商品", how="left",
    validate="many_to_one", indicator=True)
eligible = enriched["_merge"].eq("both") & enriched["有效数量"].gt(0)
pending = enriched.loc[~eligible].copy()
known = enriched.loc[eligible].copy()
known["金额"] = known["有效数量"] * known["单价"]
print(pending[["订单", "数量", "商品", "_merge"]])
\end{lstlisting}
\end{frame}

\begin{frame}[fragile]{16. 汇总：金额与覆盖范围一起报告}
\small 已知金额日报仅覆盖四笔订单，不能称为全部销售额。汇总前后金额应一致，同时报告两笔待核订单。
\medskip
\begin{lstlisting}
known_report = known.groupby(["日期", "门店"], as_index=False).agg(
    已知金额=("金额", "sum"), 可计算订单数=("订单", "size"))
print(known_report)
print("待核订单数：", len(pending))
assert len(raw_orders) == 7 and len(tidy_orders) == 6
assert len(known) == 4 and len(pending) == 2
assert known_report["已知金额"].sum() == 110
\end{lstlisting}
\end{frame}

\begin{frame}[fragile]{17. 变形：缺失单元格不能随意填零}
\small 日报透视成“日期×门店”方便比较。西店首日有数量待核订单，不能把该格写成零销售；即使已有金额的格子也可能尚未完整。
\medskip
\begin{lstlisting}
known_wide = known_report.pivot(index="日期", columns="门店", values="已知金额")
pending_wide = pending.groupby(["日期", "门店"]).size().unstack("门店")
print("已知金额（不是完整销售额）：")
print(known_wide)
print("待核订单数：")
print(pending_wide)
\end{lstlisting}
\end{frame}
\begin{frame}[fragile]{课堂练习：17. 变形：缺失单元格不能随意填零}
核实订单03数量为1、商品X是5元的水。重新清洗合并并汇总，完整金额应为125元。不要在最终报表中直接改金额。
\medskip
\begin{block}{检查思路}
先写出预期结果，再运行。参考答案见配套 Notebook 的折叠单元格。
\end{block}
\end{frame}

\section{金融时间序列}
%==============================================================

\begin{frame}[fragile]{把字符串变成日期}
	\begin{lstlisting}
df.index = pd.to_datetime(df.index)
df.index       # DatetimeIndex(['2026-07-01', ...])

pd.date_range("2026-01-01", periods=5, freq="B")
# 从1月1日起的5个工作日（B = Business day）
\end{lstlisting}
	\begin{itemize}
		\item \texttt{pd.to\_datetime()}：把字符串转成真正的日期类型。
		\item 变成日期后，可以按年、月、星期\alert{智能取值与聚合}。
		\item \texttt{date\_range} 生成规则日期序列，\texttt{freq="B"} 表示周一至周五，不排除交易所节假日。
	\end{itemize}
	\begin{exampleblock}{日期索引的威力}
		\texttt{df.loc["2026-07"]}：直接取出2026年7月的全部数据。
	\end{exampleblock}
\end{frame}

\begin{frame}[fragile]{收益率计算：pct\_change 与 shift}
	\begin{lstlisting}
df.pct_change(fill_method=None)        # 每列的日收益率：(今日-昨日)/昨日
df.shift(1)            # 整表向下错位一行 -> 昨日价格
df.diff()              # 价格变动额（今日-昨日）
\end{lstlisting}
	\begin{itemize}
		\item \texttt{pct\_change()}：把第3章的手写公式封装成了一个方法，\\
			整张表所有列\alert{一次全算}。
		\item \texttt{shift(1)}：错位是金融计算的核心技巧——
			"昨日收益""信号与次日收益配对"都靠它。
		\item 第一行没有昨日数据，结果自动为 \texttt{NaN}（缺失）。
	\end{itemize}
\end{frame}

\begin{frame}[fragile]{重采样：日频转月频}
	\begin{lstlisting}
df.resample("ME").last()    # 每月最后一个交易日的价格
df.resample("ME").mean()    # 每月平均价格
\end{lstlisting}
	\begin{itemize}
		\item \texttt{resample(频率)}：把时间序列改换频率，像"汇总报表"。
		\item 常用频率：\texttt{"D"} 日、\texttt{"W"} 周、\texttt{"ME"} 月末、\texttt{"YE"} 年末。
		\item 后面必须接聚合方式：\texttt{last/mean/sum} 等。
	\end{itemize}
	\begin{exampleblock}{典型应用}
		把日度净值换算成月度收益报告；把分钟线合成日线。
	\end{exampleblock}
\end{frame}


%==============================================================
\section{缺失值处理}
%==============================================================

\begin{frame}[fragile]{认识NaN}
	\begin{lstlisting}
import numpy as np
s = pd.Series([1500.0, np.nan, 1515.0])

s.isna()        # [False, True, False]  找出缺失
s.dropna()      # 删掉缺失的行
s.fillna(s.mean())   # 用均值填充
s.ffill()       # 用前一个有效值填充（金融常用！）
\end{lstlisting}
	\begin{itemize}
		\item \texttt{NaN}（Not a Number）表示缺失：停牌、数据未披露、\texttt{pct\_change} 的第一行。
		\item \texttt{ffill()}（前向填充）：沿用上次观测值；是否适合需依据数据含义判断。
		\item 处理前先想：缺失代表什么经济含义？乱填会引入偏差。
	\end{itemize}
\end{frame}


%==============================================================
\section{读写数据}
%==============================================================

\begin{frame}[fragile]{read\_csv：读入真实数据}
	\begin{lstlisting}
df = pd.read_csv("../data/eod_data.csv",       # 文件路径
                 index_col=0,          # 第0列作为索引
                 parse_dates=True)     # 把索引解析为日期
\end{lstlisting}
	\begin{itemize}
		\item 真实行情/财报数据多以 CSV、Excel 形式存在。
		\item 读入时顺手完成"第一列当日期索引"，后面全受益。
		\item 导出同样简单：\texttt{df.to\_csv("../outputs/result.csv")}。
		\item Excel：\texttt{pd.read\_excel("../data/data.xlsx")}（需安装 openpyxl）。
	\end{itemize}
	\begin{alertblock}{路径小贴士}
		从 \texttt{notebooks/} 运行时，\texttt{../} 表示本章目录；\\
		输入放 \texttt{data/}，生成结果放 \texttt{outputs/}。
	\end{alertblock}
\end{frame}


%==============================================================
\section{综合案例与小结}
%==============================================================

\begin{frame}[fragile]{综合案例：三股风险分析（构造数据）}
	\begin{lstlisting}
rng = np.random.default_rng(42)
dates = pd.date_range("2026-01-05", periods=250, freq="B")

# 模拟三只股票的日对数收益率，再还原为价格
rets = pd.DataFrame({
    "茅台": 0.0004 + 0.015 * rng.standard_normal(250),
    "平安": 0.0002 + 0.020 * rng.standard_normal(250),
    "招行": 0.0003 + 0.018 * rng.standard_normal(250)},
    index=dates)
prices = 100 * np.exp(rets.cumsum())
prices.tail(3)          # 查看最后3天
\end{lstlisting}
	\begin{itemize}
		\item 用 DataFrame 字典一次性构造三列，索引是简化的工作日序列。
		\item \texttt{cumsum}、\texttt{exp} 对整张表向量化生效——NumPy 功力再现。
	\end{itemize}
\end{frame}

\begin{frame}[fragile]{综合案例：收益率、风险与相关性}
	\begin{lstlisting}
rets = prices.pct_change(fill_method=None).dropna()   # 日收益率

rets.describe().round(4)              # 统计摘要
rets.mean() * 252                     # 年化平均收益
rets.std() * np.sqrt(252)             # 年化波动率
rets.corr().round(2)                  # 相关系数矩阵
\end{lstlisting}
	\begin{itemize}
		\item 一个方法作用于\textbf{所有列}：这就是 DataFrame 的威力。
		\item \textbf{相关系数矩阵}：衡量资产间联动性，组合分散化的依据。
		\item 从构造数据到风险报告，没有写一个循环。
	\end{itemize}
	\begin{exampleblock}{课堂讨论}
		相关系数接近1的两只股票，能互相分散风险吗？
	\end{exampleblock}
\end{frame}

\begin{frame}{本章小结}
	\begin{itemize}
		\item \textbf{两种结构}：Series（一列带标签）、DataFrame（一张表）；\alert{行=日期，列=资产}。
		\item \textbf{选取}：\texttt{df[列名]} 取列；\texttt{loc} 按标签、\texttt{iloc} 按位置；布尔条件筛行。
		\item \textbf{排序分组}：\texttt{sort\_values/sort\_index}；\texttt{groupby} 三步走。
		\item \textbf{时间序列}：\texttt{to\_datetime}、\texttt{pct\_change}、\texttt{shift}、\texttt{resample}。
		\item \textbf{缺失值}：\texttt{isna/dropna/fillna/ffill}。
		\item \textbf{读写}：\texttt{read\_csv(index\_col=0, parse\_dates=True)}。
		\item \textbf{金融工具箱}：\texttt{describe()}、年化波动率 $\sigma\times\sqrt{252}$、相关矩阵 \texttt{corr()}。
	\end{itemize}
\end{frame}

\begin{frame}{课后任务与下章预告}
	\begin{block}{课后任务}
		\begin{enumerate}
			\item 运行课堂 notebook，替换 \texttt{freq}、种子等参数观察输出。
			\item 对案例数据：找出每只股票的最大单日跌幅及其日期\\
			（提示：\texttt{idxmin()} 返回最小值所在的索引）。
			\item 选做：用 \texttt{resample("ME").last()} 得到月末价格，\\
			再算月度收益率。
		\end{enumerate}
	\end{block}
	\begin{exampleblock}{下章预告}
		数字算出来了，怎么"看"得清楚？\\
		下一章学习 \textbf{Matplotlib 可视化}：价格曲线、收益分布、K线之外的图表语言。
	\end{exampleblock}
	\begin{center}
		\vspace{0.5em}
		本章内容改编自教材第5章\citep{hilpisch2018python}。
	\end{center}
\end{frame}


%==============================================================
%  附录：参考文献 + 结尾页
%==============================================================
\begin{frame}{扩展阅读与练习来源}
\begin{itemize}
\item Wes McKinney, \textit{Python for Data Analysis}, 3rd ed.
\item \url{https://wesmckinney.com/book/pandas-basics}
\item 新增中文说明、合成数据与练习为本课程自行编写。
\item 阅读正文理解概念；在配套 Notebook 中运行、修改、核对。
\end{itemize}
\end{frame}

\appendix

\begin{frame}[allowframebreaks]{参考文献}
	\bibliographystyle{style/gbt7714-2005}
	\bibliography{reference.bib}
\end{frame}

\begin{frame}
	\begin{center}
		{\Huge\calligra Thanks for your attention!}
		\vspace{1cm}

		{\Huge Q \& A}
	\end{center}
\end{frame}

\end{document}
